\section{What the uniform electron gas does and does not explain}
\label{sec:ueg}

\subsection{Uniform equilibrium does not mean only zero momentum}

The UEG has a constant equilibrium density, but it can be perturbed at every
wavevector.  A scalar external potential has the form
\begin{equation}
  \delta V(\vr)=\delta V_{\vq}\ee^{\ii\vq\cdot\vr},
  \qquad
  \hat\rho_{\vq}=\sum_{\vk}c^\dagger_{\vk+\vq}c_{\vk}.
  \label{eq:rhoq}
\end{equation}
Translational invariance makes different wavevectors independent,
$\chi(\vq,\vq')=\delta_{\vq\vq'}\chi(\vq)$; it does not eliminate finite
$\vq$.  Only $\hat\rho_{\bm 0}=\hat N$ is conserved.  At finite $\vq$ the
perturbation is a density wave, and no exact identity says its quasiparticle
coupling is unrenormalized.

For elastic scattering on a spherical Fermi surface,
\begin{equation}
  q=2k_F\sin(\theta/2),
  \label{eq:qangle}
\end{equation}
so physically relevant scattering spans $0\le q\le2k_F$.  A theory that works
only at $q=0$ does not yet predict transport or pairing.

\subsection{The quantitative DFPT--many-body matching condition}

Using the scalar relations of \cref{sec:dfpt,sec:manybody}, define the
interacting irreducible polarization through a static local kernel,
\begin{equation}
  P(q,0)=\frac{\chi_s(q,0)}{1-f_{\mathrm{xc}}(q,0)\chi_s(q,0)}.
  \label{eq:pfromfxc}
\end{equation}
Then
\begin{equation}
  1-(v+f_{\mathrm{xc}})\chi_s
  = (1-vP)(1-f_{\mathrm{xc}}\chi_s).
  \label{eq:factorization}
\end{equation}
Comparing \cref{eq:gdfptscalar,eq:gphysical} shows that equality between the
two scalar vertices requires
\begin{equation}
  z\Gamma_\rho(q,0)
  \simeq \frac{1}{1-f_{\mathrm{xc}}(q,0)\chi_s(q,0)}
  = \frac{P(q,0)}{\chi_s(q,0)}.
  \label{eq:matchingcondition}
\end{equation}
This is the quantitative relationship among quasiparticle residue, vertex, and
static response relevant to the cited electron-gas comparison.  It is not the
same as setting $z\Gamma_\rho=1$, and it is not produced by the dynamic Ward
identity alone.

Recent vertex-diagrammatic Monte Carlo calculations for the electron liquid at
$r_s=1$--$5$ report that the fully interacting quasiparticle vertex and a
DFPT-like expression with a static exchange--correlation kernel agree accurately
over momenta extending to $2k_F$ \citep{Cai2025}.  The residual dependence on the
incoming fermion momentum is weak, with backscattering providing the most
demanding comparison.  This is nontrivial numerical evidence for
\cref{eq:matchingcondition} across finite momentum; it is not an exact theorem.

\subsection{Why the UEG is exceptionally simple}

The UEG contains a scalar density channel but no atomic orbitals, sublattices,
crystal-field splittings, bond angles, hybridization matrix, Umklapp mixing, or
reciprocal-space local-field matrix.  Its low-energy perturbation space is
therefore effectively one-dimensional.  Moreover, local and semilocal density
functionals are built from electron-gas information, making the UEG the most
favorable setting for a density-kernel closure.

This explains both the value and the limitation of the UEG result:

\begin{itemize}
  \item it provides a controlled baseline for the total-density channel and a
  candidate finite-$q$ prior;
  \item it cannot determine the correction to orbital, bond, hopping, spin, or
  interaction-modulation channels that do not exist in the model; and
  \item its finite-$q$ agreement must be tested, not assumed, when transferred to
  a crystal with local-field and multi-orbital structure.
\end{itemize}

\begin{keypoint}[UEG conclusion]
The UEG does not show that real materials have only $q=0$ coupling.  It shows that
when the perturbation space contains only density, the many-body problem may be
compressed into an accurate static density kernel over a surprisingly wide
momentum range.  The next-generation theory asks whether an analogous compression
exists channel by channel in solids.
\end{keypoint}
